% *- coding: UTF-8 -*-
% !TEX builder = Texlive
% !TEX prggram = xelatex

%文档设置及格式
\documentclass[UTF8,a4paper,12pt,titlepage]{ctexart}
\include{document_setting}

% 设置文件属性内容
\hypersetup{pdfauthor={吴毓明},
            pdftitle={高锰酸钾溶液制备投加一体化装置},
            pdfsubject={使用说明},
            pdfkeywords={任何问题，联系61508712@qq.com}}

% 定义页眉页脚
% 定义页眉页脚
\usepackage{fancyhdr}
\fancyhf{}
\pagestyle{fancy}
% 公共
\rhead{\zihao{6}一体化全自动加药装置/化工泵\\
       Manage Pump as Diamond} % 页眉右

% % 飞源
%\lhead{广州市花都自来水有限公司} % 页眉左
%\rhead{花都水厂应急投加装置}
%\lfoot{广州汇泽源环保科技有限公司}
%\rfoot{\thepage} % 页脚右

% 泊多科技
%\lhead{西江源水管理所} % 页眉左
%\rhead{高锰酸钾上料设备}
%\lfoot{泊多科技（广州）有限公司}
%\rfoot{\thepage} % 页脚右

% 利亨昌
% 简略版
\lhead{广州市花都自来水有限公司} % 页眉左
\rhead{花都水厂应急投加装置}
\lfoot{广州市利亨昌贸易有限公司}
\rfoot{\thepage} % 页脚右

% 详细版
% \lhead{\includegraphics{LHC-head.png}}% 页眉左
% \lhead{\includegraphics[height=0.5cm]{LHC.jpg}} % 页眉左
% \chead{\zihao{4}广州市利亨昌贸易有限公司}
% \chead{\zihao{5}广州市利亨昌贸易有限公司}
% \lfoot{\zihao{6}广州市荔湾区东联路2号工业园\\
%                电话：020-8430 1397 传真：020-84301170 \\ 
%                电邮：gzlhcmy@163.com 网址：www.lhcmy.com}
% \cfoot{\thepage} % 页脚右
% \cfoot{\zihao{5}\thepage} % 页脚中
% \rfoot{\zihao{6}美国帕斯菲达PULSAFEEDER 产品库存商\\Paddle-flow流量监控及加药装置} % 页脚右
% \rfoot{\zihao{6}美国帕斯菲达PULSAFEEDER 产品库存商\\Paddle-flow流量监控及加药装置} % 页脚右
% \rfoot{\thepage} % 页脚右

% paddle-flow
%\rhead{\includegraphics{paddle-flow.png}}

% 厦门飞华
%\lhead{宁夏宁东水务有限责任公司} % 页眉左
%\rhead{高锰酸钾制备投加一体化装置}
%\lfoot{厦门飞华环保有限公司}
%\rfoot{\thepage} % 页脚右

% 王远望王老师
%\lfoot{厦门创洁数智环保水处理科技有限公司}
%\rfoot{\thepage} % 页脚右

% 杨海涛
%\lhead{宁夏海原水厂}
%\rhead{次氯酸钠加药装置}
%\lfoot{宁夏智诚博创自动化科技有限公司}
%\rfoot{\thepage} % 页脚右

\renewcommand{\headrulewidth}{0.4pt} % 页眉分隔线
\renewcommand{\footrulewidth}{0.4pt} % 页脚分隔线


% 表格格式定义
\include{table_format}

% 图片
\include{picture_configuration}

% 特殊符号(例如千分号等)
\usepackage{textcomp}

% 流程图
\include{flow_chart_preamble}

% 开始正文
\begin{document}

% 制作封面
\title{
   \zihao{1} 广州市花都自来水有限公司\\[2.5mm]
   \zihao{1} 花都水厂应急投加装置\\
   \zihao{1} 综合技术服务方案\\
   {\noindent}\rule{16cm}{1pt}\\[3mm]
   \zihao{2} 广州市利亨昌贸易有限公司\\
   % \zihao{4} 吴毓明\\
   %\zihao{4} Wym\\
   %\zihao{4} 2022年12月27日
   \begin{figure}[h]
      \centering
      \includegraphics[height=10cm]{g01.PNG}
   \end{figure}
}

\maketitle % 封面



\tableofcontents % 目录

\newpage % 换页

%\part{系统组成}
%ctex的\chapter{sdsg}是不能用的

\section{项目技术要求及实施目标}
    \subsection{项目技术要求}
        
%高锰酸钾制备装置采用一体化成套结构，由螺旋给料器、加热器、电磁阀、震动器、
%物位开关、超声波液位计、粉剂料斗、搅拌器、储存箱、流量计、喷射混合器、真空吸料机
%及内部管道阀门和 PLC 控制柜等组成。
%
%\par装置制备能力1000L/h，浓度：0.5$\sim$5.0\%可调。
%  
%制备装置的效果图如下所示：\\
%
%\begin{figure}[h]
%   \centering
%   \includegraphics[height=11cm]{g06.PNG}
%   \caption{一体化全自动高锰酸钾制备装置效果图}
%\end{figure}
%
%\par上图供参考，实物布局以最终设计为准。

本项目旨在为花都水厂提供一套高标准、全自动、高度集成的应急投加装置。
设备须采用“干粉上料、快速溶解、数字泵精准投加”三大模块分体可移动设计。
核心硬性技术性能指标：
\begin{itemize}
    \item 投加流量误差$\le \pm 1 \%$。
    \item 系统最大投加量$\ge 1500L/h$。
    \item 具备连续72小时无故障运行能力。
    \item 且所有涉药部件必须采用304不锈钢（经酸洗钝化）或PP/UPVC等防腐材质。
\end{itemize}


    \subsection{实施目标}
        \begin{itemize}
    \item 质量目标
        \begin{itemize}
            \item 设备制造与验收合格率100\%。
            \item 材质复检合格率100\%。
            \item 连续72小时试运行零故障。
        \end{itemize}
    \item 工期目标
        \begin{itemize}
            \item 严格遵守合同，确保在收到通知后45天内完成交货、现场安装指导及全工况联动调试。
        \end{itemize}
    \item 服务目标
        \begin{itemize}
            \item 提供零死角的现场培训，确保水厂操作人员能独立完成启停、参数调节及日常维护。
        \end{itemize}
\end{itemize}





\section{生产制造方案}
    \subsection{生产准备与原材料管控}
        \subsubsection{原材料采购与复检}
        严格按照图纸采购304不锈钢（板材、型材）、PP板材、UPVC管道/管件、
        耐腐蚀密封件（耐PAC、PAM、高锰酸钾、碳酸钠腐蚀）。
        所有进厂材料100\%核对材料材质。

        \subsubsection{下料与预处理}
        采用激光切割机进行不锈钢板材下料，
        保证切割精度；
        所有不锈钢构件在焊接前进行严格的酸洗钝化处理，
        确保防腐性能，
        严禁使用201不锈钢或碳钢喷漆件。

    \subsection{核心模块精密制造工艺}
        \subsubsection{干粉上料小车}
        干粉上料小车 由 真空上料系统、干粉料仓、管路系统组成，各系统特点如下：
        \begin{itemize}
            \item 真空上料系统：\\
                配置自带过滤器及反吹清灰功能的真空上料系统，适配吸吨包、小包干粉、罐装固体粉料。
            \item 干粉料仓：\\
                制作304不锈钢200L储料仓。
                内部集成“机械破拱+流化气嘴”双防架桥系统。
                顶部设计带料位高低报警的密闭投料口及吸尘接口，
                确保应急投料无大量粉尘。
            \item 管道系统：\\
                独立进水口及吹扫接口，
                停机时通入压缩空气吹扫管路，
                防止干粉板结堵管。
        \end{itemize}

        \subsubsection{快速溶解小车}
        \begin{itemize}
            \item 罐体：\\
                溶解罐容积1000L，
                不锈钢304或PP等腐蚀材质，
                内部精密分为预混区、溶解区、熟化区。
            \item 搅拌系统：\\
                采用“双层高速混合+底部防沉淀”搅拌器。
                低转速熟化、高速预湿，彻底解决PAM结“鱼眼”问题。
                搅拌轴及叶轮进行动平衡测试，确保运行平稳。
                搅拌机及叶轮成套不锈钢304材质。
            \item 流体控制：\\
                配置自动补水系统（液位开关+流量计双控），
                实现按浓度自动配水或纯手动补水。
                罐底设全排空排污阀，确保清洗无死角。
                粉体物料投加与补水流量互相配合，
                确保制备溶液浓度稳定。
        \end{itemize}

        \subsubsection{数字泵精准投加小车}
           数字泵精准投加小车 包含 数字泵、阻尼器、背压阀、泄压阀（安全阀）、隔膜压力表、加药流量计等组件， 
           具体说明如下：
        \begin{itemize}
            \item 数字泵：配置500L/h+1000L/h 2台数字计量泵。两既可以按一用一备使用，也可以双泵同投。
                数字泵具备 冲程 + 变频双重调节功能，并根据双脉冲流量计反馈闭环调整流量。
            \item 阻尼器：
            每台泵配置FPM隔膜气囊式脉动阻尼器，平滑计量泵输出流量脉动。 PVC材质壳体，FPM隔膜，带充气压力表。
            \item 背压阀：
            每台泵配置背压阀，PVC材质阀体，PTFE/氟橡胶复合隔膜。 背压压力0-3bar。
            \item 泄压阀（安全阀）：
            每台泵配置泄压阀，系统过压时保护系统。 PVC材质阀体，氟橡胶隔膜，背压压力0-6bar。
            \item 隔膜压力表：
            每台泵配置隔膜压力表。 PVC材质过液端，FPM隔膜，内充硅油作为传递压力的中间液。 适用于测量输送腐蚀性介质管道的压力。
            \item 加药流量计：
            每台计量泵均配置1个流量计，检测投加流量。 
            测量流量范围0-1000L/h，脉冲式或齿轮式。
            流量反馈信号4-20mA或RS-485可选。
        \end{itemize}

    \subsection{ 电气与自控系统集成}
        电气与自控系统集成 供货标准如下：
        \begin{itemize}
            \item 西门子品牌PLC，同品牌系列输入输出模块。
            \item 西门子或施耐德品牌变频器。
            \item 昆仑通泰或威纶通触摸屏。
            \item 施耐德电气元件。
            \item 电气设备具有过载、短路、缺相等多重保护。
            \item 柜内布线整齐，线号标识清晰，接地可靠。
            \item 预留硬接线干接点、4-20mA模拟量及Modbus-RTU通讯接口。
            \item PLC程序无加密、无程序锁、无硬件绑定、无付费解锁，完全开放控制。
        \end{itemize}

\section{全面质量控制方案}
    \subsection{质量检验体系}
    实施“三检制”（自检、互检、专检），建立从原材料到成品的全生命周期质量追溯档案。

    \newpage

    \subsection{关键工序质控点（重点控制项目）}
    \vspace{0.5cm}

    \begin{table}[htbp]
        \centering
        \renewcommand{\arraystretch}{1.5} % 增加行高，使表格更美观
        \begin{tabular}{>{\centering\arraybackslash}p{2.5cm} >{\raggedright\arraybackslash}p{7cm} >{\raggedright\arraybackslash}p{4.5cm}}
            \toprule
            \textbf{质控点} & \textbf{检验标准与措施} & \textbf{检验方法/工具} \\
            \midrule
            \textbf{材质核对} & 严禁201不锈钢/碳钢喷漆冒充。供方提供板材、管材材质证明。 & 化学方法复检、游标卡尺测厚 \\
            \midrule
            \textbf{焊接与密封} & 焊缝平整，无气孔；管路无跑冒滴漏；罐体实现全排空，无积液死角。 & 满水保压试验（24h）、目视检查 \\
            \midrule
            \textbf{防腐处理} & 所有不锈钢构件（含机架）100\%经酸洗钝化处理。 & 表面钝化膜检测（蓝点测试） \\
            \midrule
            \textbf{装配与结构} & 整机低重心设计，机架配套配重，带刹车重型万向轮。突出部件防护无尖锐棱角。 & 推行测试、防倾倒及刹车测试 \\
            \midrule
            \textbf{电气防护} & 户外电气元件防护等级\(\ge\)IP65；独立电控柜；线缆接头防水。 & 绝缘电阻测试、接电阻仪测量 \\
            \midrule
            \textbf{自控仪表} & 变送器、流量计精度满足整机要求；支持U盘导出不少于90天运行及报警记录。 & 信号发生器模拟测试、U盘读取验证 \\
            \midrule
            \textbf{投加精度} & 数字泵“冲程+变频”双重调节，双脉冲流量计闭环校正。 & 标定柱容积测量、秒表计时 \\
            \bottomrule
        \end{tabular}
    \end{table}

    \subsection{出厂前测试}
    出厂前测试 包含 空载试运行、负荷试运行两部分。
        \begin{itemize}
            \item 空载试运行：\\
                检验各电机转向、触摸屏界面交互、
                报警逻辑（低液位停数字泵、高液位停进水、数字泵故障报警）。
            \item 负荷试运行：\\
                进行闭环校准，验证数字泵在变频+冲程双重调节下的精确度，确保出厂前流量误差$\le \pm1\%$。
        \end{itemize}

\newpage
\section{项目进度计划}
    为满足“45天内到货并安装调试完成”的要求，我司采用并行工程管理，制定如下详细节点计划：

    \vspace{0.5cm}

    \begin{table}[htbp]
        \centering
        \renewcommand{\arraystretch}{1.6} % 增加行高，使表格更美观
        \begin{tabular}{>{\centering\arraybackslash}p{2.2cm} >{\centering\arraybackslash}p{2.2cm} >{\raggedright\arraybackslash}p{7.5cm} >{\raggedright\arraybackslash}p{3cm}}
            \toprule
            \textbf{阶段} & \textbf{时间节点} & \textbf{关键工作任务分解} & \textbf{交付物/里程碑} \\
            \midrule
            \textbf{一、设计与备料} & 第1-5天 & 召开项目启动会，下发BOM清单；非标件图纸深化；原材采购进厂及材质复检。 & 项目排产计划、材质证明书 \\
            \midrule \textbf{二、生产与组装} & 第6-15天 & 机架焊接打磨；料仓与溶解罐制作（含酸洗钝化）；外购件（泵、仪表）检测。 & 半成品完工报验单 \\
            \midrule
            \textbf{三、电气集成与总装} & 第16-22天 & 独立电控柜组装及PLC编程；三车机械总装；管路连接与快接安装。 & 机械/电气装配完工单 \\
            \midrule
            \textbf{四、厂内测试(FAT)} & 第23-27天 & 空载测试；带水联动测试；数字泵精度标定；包装前质检。 & FAT测试报告、合格证 \\
            \midrule
            \textbf{五、包装与发运} & 第28-30天 & 模块化拆解，专业防震防雨包装；安排专车物流发运至花都水厂。 & 发货清单、物流单号 \\
            \midrule
            \textbf{六、现场卸货与就位} & 第31-32天 & 设备到场清点；叉车卸货；单人推行至指定位置；锁死万向轮。 & 签收确认单 \\
            \midrule
            \textbf{七、安装与接线指导} & 第33-37天 & 我方工程师现场指导：水管/药管快接对接；电源线接入；联动信号（干接点/4-20mA/RTU）接线。 & 现场安装记录表 \\
            \midrule
            \textbf{八、系统联调与试运行} & 第38-42天 & 各模块独立启停测试；全流程“上料-溶解-投加”联动调试；\textbf{连续72小时满负荷（\(\ge\)1500L/h）无故障试运行}。 & 调试报告、72h运行曲线记录 \\
            \midrule
            \textbf{九、培训与交付验收} & 第43-45天 & 开展理论+实操培训；移交全套竣工资料；配合甲方进行终验。 & 培训签到表、验收单 \\
            \bottomrule
        \end{tabular}
    \end{table}

\section{装卸、安装与调试详细方案}
    \subsection{装卸与转运方案}
        \begin{itemize}
            \item 包装保护：\\
            设备采用木托盘底座，外层用防震气泡膜及防雨缠绕膜包裹，尖锐部位加装防撞护角。
            \item 卸货与流转：\\
            设备到达现场后，
            使用叉车平稳卸至地面。
            因整机采用低重心设计及重型万向轮，无需吊车，
            操作人员单人即可将三台小车平稳推入水厂投药间内。
            就位后立即踩下万向轮刹车装置。
        \end{itemize}

    \subsection{现场安装指导步骤}
        \begin{itemize}
            \item 管线快接：\\
            指导运营人员使用统一标准快接接口连接外部水源、干粉送料管、出药管。
            无需使用胶水临时粘接，无需专用工具。
            不用的备用管口立即装上密封堵头。
            \item 电源接入：\\
            使用设备自带的防水耐油耐腐蚀电源线，接入标准快速接头即可。
            \item 联动接线：\\
            若需厂区联控，接入外部4-20mA进水流量信号及Modbus-RTU通讯线；
            若仅需独立运行，断开联动线缆即可，各模块互不影响。
        \end{itemize}

    \subsection{深度调试方案}
        \begin{itemize}
            \item 单机功能验证：\\
            验证干粉车真空吸料及反吹清灰；
            验证溶解车自动补水逻辑及低液位停数字泵、高液位停进水；
            验证投加车数字泵的手动/自动切换。
            \item 精度标定：\\
            使用标定柱，
            在全量程（0-1500L/h）内多点取样。
            利用数字计量泵的冲程调节及变频调节，结合双脉冲流量计的闭环比对，确保投加流量误差严格控制在$\pm 1\%$以内。
            \item 全工况72小时试运行：
            \begin{itemize}
                \item 模拟外部进水流量4-20mA信号变化，测试系统按比例精准投加。
                \item 满负荷运行，观察管路无跑冒滴漏，罐体无固体粉料结块，粉料系统无扬尘。
                \item 满测试故障模拟，验证报警就地提醒及自动停机保护是否生效。
                \item 验证运行记录本地存储与U盘导出功能（$\ge 90$天记录）。
            \end{itemize}
        \end{itemize}

\section{技术培训与售后交付}
    \subsection{“手把手”现场技术培训}
    \begin{itemize}
        \item 理论培训（半天）\\
        讲解工艺流程图、模块化设计理念、各部件功能及安全操作规范（如遇紧急情况一键急停）。
        \item 实操培训（一天）\\
            \begin{itemize}
                \item 启停操作：
                独立启停、联动启停操作，手动溶液制备操作，手动投加流量设定。
                \item 参数设置：
                触摸屏上设定溶液制备量L/h、制备浓度设定、查看瞬时投加流量、溶液制备量。
                \item 维护保养：
                Y型过滤器快速拆洗、管路的日常吹扫防堵、各设备齿轮箱油位检查、泵隔膜及阀组件更换。
                \item 故障排查：
                常见报警代码消除、传感器清零与维护。
            \end{itemize}
        \item 考核闭环\\
        培训结束后，对运营人员进行现场抽考，确保每人能独立完成日常操作。
    \end{itemize}

    \subsection{交付资料清单}
    随机 交付资料清单 如下：
        \begin{itemize}
            \item 设备总装图、电气原理图、管路图（CAD及PDF版）。
            \item 设备原厂合格证。
            \item 全套材质检测报告（含板材、管材光谱复检报告）。
            \item 电气检测报告（绝缘、接地测试）。
            \item 仪器仪表出厂校准证书。
            \item 出厂测试报告及现场72小时连续运行调试记录。
            \item 《设备操作与维护手册》。
            \item 自控系统PLC控制程序源程序。
        \end{itemize}

\section{售后保障}
提供质保期内免费维修及终身技术指导。
建立专属客户服务档案，提供7×24小时电话响应，
如需现场支持，承诺在接到通知后合同约定的最快时间内抵达花都水厂现场。

\end{document}
